\documentclass[aps,prl,reprint,superscriptaddress,nofootinbib,floatfix]{revtex4-2}

\usepackage{amsmath,amssymb,bm}
\usepackage{graphicx}
\usepackage{xcolor}
\usepackage[colorlinks=true,citecolor=blue,linkcolor=blue,urlcolor=blue]{hyperref}

\begin{document}

\title{Geometry-Encoded Hidden Orientational Memory}

\author{Vladimir A. Baulin}
\affiliation{Synthetix Institute, Amsterdam, The Netherlands}
\email{vbaulin@synthetix.institute}

\date{}

\begin{abstract}
Memory in disordered soft matter usually requires structural rearrangement or
a persistent external template. Whether a material can instead retain a
microscopic pattern in mobile internal coordinates remains unresolved. Here
we study elongated Brownian colloids whose centres are caged at a liquid
interface while their axes rotate. Capillary quadrupoles register each pair
against its bond direction. On a disordered neighbour graph, these local
conditions compete and prevent a common director. Simulations show that two
replicas descended from one prepared configuration retain an angular overlap
of $0.462\pm0.010$ after 625 rotational-diffusion times. A transiently
written state retains $0.441\pm0.021$ after the field is removed. A matched
system without the capillary channel loses both records within a few
rotational times. Spatial correlations end within two particle spacings, and
finite-size tests find no equilibrium spin-glass crossing over the sizes
examined. Retention is largest near the disorder at which global nematic order
collapses. The result identifies a template-free route by which fixed bond
geometry can encode long-lived orientational memory.
\end{abstract}

\maketitle

A material memory links a present response to a past perturbation. In soft
matter, that link is commonly carried by rearranged particles, hysteretic
elements, defects, or fixed alignment fields \cite{keim2019}. Cyclically
sheared suspensions, for example, reorganize until their trajectories encode
the training amplitude \cite{keim2011,paulsen2014}. Quenched-disorder nematics
instead retain local easy axes \cite{rotunno2005}. Both mechanisms store the
record outside a freely evolving orientation. Rotational arrest provides a
different contrast. Colloidal ellipsoids exhibit orientational and
liquid-glass arrest \cite{zheng2011,roller2021}. Slow angles, however, need not
preserve a specific preparation after its field is removed.
We therefore use a stricter criterion. A preparation must remain identifiable
after its conjugate field is absent. Its overlap must also exceed a matched
system in which the proposed storage interaction is removed. This separates a
memory from slow relaxation measured over one finite window.

Here we ask whether pair interactions can store such a preparation while the
centres and external fields remain unchanged. The key distinction is between
relative-angle and bond-frame coupling. Nematic alignment constrains
\(\theta_i-\theta_j\). A bond-frame interaction instead constrains
\(\theta_i+\theta_j-2\phi_{ij}\), where \(\phi_{ij}\) is the bond direction.
Several differently oriented bonds cannot generally satisfy both conditions
with one director. Closed paths then divide the angular landscape into
graph-dependent compromises.
Each bond can still be locally well satisfied. The frustration appears when
those choices are continued around a loop and fail to close. A prepared state
can then select one compromise without producing a laboratory-frame nematic
axis. The memory is hidden from a bulk polarizer but remains accessible to a
particle-resolved overlap.

Interfacial ellipsoids provide this interaction directly. Their leading
contact-line deformation is quadrupolar,
\(u_i\propto\Delta u_{\max}(R/r)^2\cos2(\phi-\theta_i)\). Meniscus overlap
gives
\(U^{\rm cap}_{ij}=-A_4(r_{ij})\cos2(\theta_i+\theta_j-2\phi_{ij})\), with
\(A_4=3\pi\gamma(\Delta u_{\max})^2(R/r_{ij})^4\)
\cite{loudet2005,luo2019}. Here \(\gamma\) is the interfacial tension and
\(R\) is a particle length scale. This static deformation supplies a
persistent pair torque without a one-particle easy axis.
For \(A_4>0\), an isolated pair minimizes its energy when
\(\theta_i+\theta_j=2\phi_{ij}+m\pi\). A particle in a dense layer belongs to
several such pairs. Disorder changes their bond angles, so one orientation
cannot minimize every adjacent interaction. This realizes geometric
frustration using the particle positions themselves.

\begin{figure*}[t]
 \includegraphics[width=\textwidth]{capillary_prl_figures/fig1_capillary_realization.pdf}
\caption{\label{fig:realization}
 Physical realization. (a) Quadrupolar menisci couple two ellipsoids to their
 bond direction. (b) A caged disordered monolayer at
 \((J,g)=(4,5)k_{\rm B}T\); grey lines show the persistent neighbour graph.
 (c) The far-field coupling for \(\gamma=10\,\mathrm{mN\,m^{-1}}\).
 Nanometre menisci place separated particles in the Brownian energy range.}
\end{figure*}

Figure~\ref{fig:realization} makes the physical separation of roles explicit.
The meniscus in panel (a) generates a pair constraint, not a local easy axis.
The grey network in panel (b) records which particles interact, while the rod
angles remain dynamical. The irregular orientations are therefore not an
imposed domain pattern. They are compromises selected by the fixed bond
geometry. Panel (c) identifies the required interaction window. A quadrupole
below thermal energy decorrelates, whereas a much stronger quadrupole would
aggregate or immobilize the particles. Nanometre contact-line amplitudes place
the torque between these limits without patterned grooves.

We model \(N\) apolar Brownian rotors on a fixed neighbour graph. Their energy
is
\begin{align}
 U={}&-\sum_{(ij)}J w^-_{ij}\cos2(\theta_i-\theta_j)\nonumber\\
 &-\sum_{(ij)}g w^+_{ij}
 \cos2(\theta_i+\theta_j-2\phi_{ij}).
 \label{eq:model}
\end{align}
The first term is the leading short-range apolar alignment, with energy
\(J\) and weight \(w^-_{ij}\). It can arise from steric, near-field capillary,
or other shape-dependent interactions. The second term is the quadrupolar
channel, with \(g=A_4(r_0)\) and
\(w^+_{ij}=\min[6,(r_0/r_{ij})^4]\). The median nearest-neighbour distance is
\(r_0\). Particle centres are displaced from a triangular cage by a Gaussian
width \(\sigma\); no orientational blocks are imposed. The Supplemental
Material specifies both weights, graph construction, and simulation schedule.
The graph remains fixed during each angular trajectory. Equation~\eqref{eq:model}
therefore tests the hierarchy in which centre motion is caged while rotation
remains Brownian. Temporary writing fields are absent during every reported
retention interval.

\begin{figure*}[t]
 \includegraphics[width=\textwidth]{capillary_prl_figures/capillary_regime_diagram.pdf}\\[-1.5mm]
\caption{\label{fig:map}
 Finite-size dynamical taxonomy from 441 coupling cells on three disordered
 \(N=400\) graphs. Clustering uses only the measured observables. (a) Four
 response classes; curves are nearest-centroid boundaries. (b) Classification
 margin. (c) Cut at \(g=4.875k_{\rm B}T\). The star marks the independently
 tested point \((J,g)=(4,5)k_{\rm B}T\).}
\end{figure*}

The angular density obeys the rotational Smoluchowski equation. Its It\^o
Langevin form is integrated with rotational diffusivity \(D_r\). Global,
relative-pair, and bond-frame order are
\begin{align}
 S&=\left|N^{-1}\sum_i e^{2i\theta_i}\right|,\nonumber\\
 C_2&=\frac{\sum_{(ij)}w^-_{ij}\cos2(\theta_i-\theta_j)}
 {\sum_{(ij)}w^-_{ij}},\nonumber\\
 G_2&=\frac{\sum_{(ij)}w^+_{ij}
 \cos2(\theta_i+\theta_j-2\phi_{ij})}{\sum_{(ij)}w^+_{ij}}.
 \label{eq:order}
\end{align}
Equation~\eqref{eq:order} separates common-director order from relative-angle
and bond-frame organization. The finite-window statistic
\(q_{\rm EA}=N^{-1}\sum_i|\langle e^{2i\theta_i}\rangle_t|^2\) measures
persistence, not memory by itself. Memory is the post-release overlap with a
known preparation. We use
\(Q_{\rm split}=N^{-1}\sum_i\cos2(\theta_i^{a}-\theta_i^{b})\) for descendants
with independent noise. We use
\(Q_{\rm target}=N^{-1}\sum_i\cos2(\theta_i-\theta_i^{\rm target})\) for a
written state.
The first protocol asks whether two descendants inherit the same basin. The
second asks whether a specified preparation remains readable. We call the
response hidden orientational memory only when \(S\) is small, local
correlations are finite, and a post-release overlap survives. Equilibrium
replica tests are evaluated separately.

Figure~\ref{fig:map} resolves four finite-size response classes in the
observable vector \((S,C_2,G_2,q_{\rm EA},C_{\rm window})\). The coupling
coordinates are withheld from clustering. All 27 graph bootstraps select four
classes. Dropping either persistence coordinate still selects four, with
adjusted Rand indices 0.821 and 0.805. The map is therefore a dynamical
taxonomy, not a thermodynamic phase diagram.
The four classes distinguish a Brownian crossover, relative alignment,
capillary-frame persistence, and a mixed state. Both pair harmonics remain
finite in the mixed state while \(S\) is small. This population contains 169
of 441 cells. Its centroid is
\((S,C_2,G_2,q_{\rm EA})=(0.112,0.638,0.414,0.628)\).
The marked point lies inside this population and was selected before the
larger graphs and long trajectories were analyzed.

The arrangement of the four classes follows directly from the competing
terms in Eq.~\eqref{eq:model}. Weak \(J\) and \(g\) give the Brownian
crossover. Large \(J\) with weak \(g\) selects relative alignment and a common
director. Large \(g\) with weak \(J\) registers orientations to the bond
frame but weakens neighbour-parallel order. The mixed class occupies the
region where neither interaction eliminates the other. In Fig.~\ref{fig:map}(b),
the dark bands are cells nearly equidistant from two centroids. They measure
classification ambiguity, not thermodynamic coexistence. The cut in panel (c)
shows how increasing \(J\) builds \(C_2\) while \(G_2\) remains finite. The
marked state lies beyond this crossover, where both constraints act but
global order remains weak.

The marked hidden-memory point was fixed before the long runs. Five
\(N=1024\) graphs give
\((S,C_2,G_2,q_{\rm EA})=(0.067,0.627,0.447,0.658)\). Graph standard
deviations are below 0.011. Removing the capillary term sets \(G_2\) to zero and
reduces \(q_{\rm EA}\) to \(0.022\), even though \(S\) and \(C_2\) increase.
On a regular graph the same couplings instead give \(S=0.945\). Thus the
bond-frame term suppresses a common director while preserving local angular
constraints. Varying positional disorder alone confirms the mechanism:
retention peaks near \(\sigma=0.11a\), where nematic order first collapses
(Supplemental Figs.~S3 and S4).
Across the wider disorder scan, \(S\) falls thirteenfold, from 0.916 to 0.069,
whereas \(q_{\rm EA}\) changes by only a factor of 1.3. Above the retention
maximum, the connected written overlap decreases from \(0.559\pm0.022\) at
\(\sigma=0.11a\) to \(0.386\pm0.026\) at \(0.28a\). Disorder first removes
the common director, but stronger disorder also weakens the compatible local
network that carries the record.

\begin{figure*}[t]
 \includegraphics[width=\textwidth]{capillary_prl_figures/fig3_capillary_memory_dynamics.pdf}
\caption{\label{fig:dynamics}
 Angular memory at \(N=1024\) and \((J,g)=(4,5)k_{\rm B}T\); bands are
 graph standard deviations. (a) Descendants of one state retain overlap for
 \(D_rt=625\). (b) A written state remains after its field is removed.
 (c) Older quenched states relax more slowly. The \(g=0\) controls erase both
 records within a few rotational times.}
\end{figure*}

The dynamical tests establish memory beyond static graph registration.
Figure~\ref{fig:dynamics}(a) gives
\(Q_{\rm split}=0.462\pm0.010\) at \(D_rt=625\); the matched control falls
below 0.2 by \(D_rt=2.7\pm0.3\). A temporary field produces
\(Q_{\rm target}=0.863\pm0.009\), followed by
\(0.441\pm0.021\) after release [Fig.~\ref{fig:dynamics}(b)]. The control
erases the same target within \(4.3\pm0.3\) rotational times. The aging
increment is positive on every tested graph. These observations identify a
history-conditioned basin rather than a slowly relaxing interface. The
estimated interface relaxation time is at least three orders of magnitude
shorter (Supplemental Material).
The two-time correlation also grows with waiting time. At fixed lag
\(D_r\Delta t=30\), the paired increase is
\(0.0253\pm0.0019\), and it is positive on every graph. This finite-size aging
is consistent with slow exploration of the same multibasin landscape. It is
not used as evidence for an equilibrium glass transition.

The three panels test complementary consequences of the same landscape.
Panel (a) asks whether two noisy descendants retain a common basin without
specifying which basin it is. Panel (b) asks whether a chosen pattern can be
written and recovered after its field disappears. Panel (c) asks whether the
relaxation law depends on the age of the prepared state. Agreement of (a) and
(b) separates retained preparation from an arbitrary nonzero autocorrelation.
The aging trend shows that the basin population evolves during the quench. It
does not supply a separate memory definition.

\begin{figure*}[t]
 \includegraphics[width=\textwidth]{capillary_prl_figures/fig4_activated_memory_results.pdf}
\caption{\label{fig:activated}
 Coupling-dependent retention along
 \((J,g)=\lambda(4,5)k_{\rm B}T\), averaged over five \(N=576\) graphs.
 (a) Split-replica overlap. (b) Endpoint overlaps for descendants, written
 states, and \(g=0\) controls. (c) Finite-window \(q_{\rm EA}\). Its window
 dependence distinguishes persistence from the retained-memory readout in
 (b).}
\end{figure*}

Figure~\ref{fig:activated} resolves a coupling-controlled onset. At
\(D_rT_{\rm obs}=625\), descendant overlap rises from values consistent with
zero below \(\lambda=0.6\) to \(0.405\pm0.008\) at 0.9 and
\(0.691\pm0.005\) at 1.4. Written-state retention rises in parallel, whereas
the \(g=0\) branch remains zero. Panel (c) also shows why \(q_{\rm EA}\) is a
diagnostic rather than the memory readout.
At \(\lambda=0.30\), \(q_{\rm EA}\) decays to zero as the observation window
grows. At 1.40 it remains near 0.78 over the longest window. In contrast,
Fig.~\ref{fig:activated}(b) compares every coupling at the same read time and
against the same null channel. The onset is therefore not created by fitting
a decay law or by integrating up to a finite-window ceiling.

Panel (a) also shows that the onset is broad rather than singular. Increasing
\(\lambda\) extends the time over which overlap survives and raises its
endpoint. Panel (b) reads every trajectory at one common physical time, so
the increasing branch cannot result from a longer integration window. The
written and split-replica observables rise together, although they start from
different preparations. This links the onset to autonomous basin retention
rather than to the writing field. Panel (c) supplies the necessary contrast:
finite-window persistence can be large before its long-window limit is known.

Finite-size scaling separates this memory from equilibrium glass order.
Between \(N=144\) and 2304, \(S\propto N^{-0.479}\) and independently
equilibrated replica overlap scales as \(N^{-0.466}\). Local order remains
finite, while \(\xi_L/L\) has no common crossing. Spatial correlations reach
the noise floor by \(r/a=1.73\), yet descendants retain overlap for hundreds
of rotational times. The stored variable is therefore local basin identity,
conditioned on a common preparation.
The overlap susceptibility saturates with size, and the largest measured
second-moment correlation length is \(1.49a\). Independent preparations do
not select one macroscopic state. Descendants of the same preparation remain
related because they begin in the same local basin. This is genealogical
memory rather than equilibrium replica order.

A strong-locking intervention identifies the origin of those basins. The
apolar reduction assigns each local region a binary state and each loop the
gauge-invariant product \(\mathcal W_\gamma\) of its coupling signs. Removing
loops with \(\mathcal W_\gamma=-1\), while preserving the graph and coupling
magnitudes, removes \(5.58\) one-flip-stable states on average. Closed-path
incompatibility creates additional choices; preparation selects among them.
The full intervention and target-resolved controls are in the Supplemental
Material.
The graph intervention changes topology of the constraints without changing
their magnitudes. Across 24 graphs, the excess stable-state count has a
bootstrap 95\% interval of 2.67--8.83. Individual written patterns can improve
or worsen when loop mismatch is removed. Thus incompatible loops create
alternatives; they do not guarantee that every target will be retained.

The same landscape records operation order. Two opposed write fields applied
as \(AB\) or \(BA\) leave a signed post-release difference. At \(N=256\),
overlapping supports give \(M_{\rm ord}=0.5781\pm0.0030\), whereas separated
supports give \(0.0015\pm0.0010\). A nonzero response during writing is not
sufficient: the two histories must remain in different basins after both
fields are removed. The protocol and reduced-model control appear in
Supplemental Fig.~S9.
Equivalently, field-free relaxation defines which driven states become
indistinguishable. Pulse order is stored only when the \(AB\) and \(BA\)
histories relax to different readout classes. This makes the signed
post-release overlap a material order parameter for sequence, rather than a
measure of transient noncommutativity during writing.

The support controls reveal why overlap matters. Partitioned fields address
disjoint rotors, so reversing their order does not reverse a shared local
decision. Their driven trajectories may differ, but field-free relaxation
maps them to the same readout class. Contested fields instead favor opposite
orientations on the same rotors. The last pulse places those coordinates on
opposite sides of a basin boundary in the two histories. Coupling transmits
that choice beyond the written set and stabilizes it after removal. Sequence
memory thus requires nonuniform overlap between operations in this model.
It is not determined by the instantaneous noncommutativity of the driven
dynamics alone; the difference must survive the relaxation map.

The fixed-graph approximation requires
\(\tau_{\rm int}\ll D_r^{-1}\ll\tau_{\rm cage}\). Charged rods already show
caged translation with rotational mobility \cite{besseling2026}, and Brownian
colloidal pivots demonstrate accessible angular motion \cite{melio2026}. A
candidate realization is a dense air--water monolayer of polymer ellipsoids.
Electrostatic repulsion and crowding maintain the cage, while roughness and
wetting tune capillary attraction \cite{trevenen2023,rahman2026}. At
\(\gamma=10\,\mathrm{mN\,m^{-1}}\) and \(r/R=3\), the target
\(A_4=1\)--\(5k_{\rm B}T\) requires
\(\Delta u_{\max}=1.9\)--\(4.2\,\mathrm{nm}\).
The condition \(\tau_{\rm int}\ll D_r^{-1}\) makes the interface follow the
rotor quasistatically. The condition \(D_r^{-1}\ll\tau_{\rm cage}\) allows
the angles to explore their landscape before the
neighbour graph changes. Smooth ellipsoids interact much more strongly, so
the experiment requires an intermediate quadrupole rather than maximal
capillary attraction.

The same interface provides three experimentally distinct controls. A
polarized field writes \(\theta_i\) without moving the centres. Deviatoric
surface curvature supplies a mechanical write torque and vanishes when the
surface is flattened. A spatially uniform capillary wave instead modulates
\(g\), changing the escape barrier without encoding a target. Waves that move
the centres change \(\sigma\) and therefore traverse the retention maximum
near \(0.11a\). These operations predict separate write, gate, and erase
responses that can be distinguished by tracking both positions and angles.

Video microscopy provides every model observable from \(\bm r_i(t)\) and
\(\theta_i(t)\). A two-particle measurement first calibrates \(J\), \(g\),
and the \(r^{-4}\) range. The collective test then perturbs the layer,
removes the perturbation, and compares post-release overlap with a
low-quadrupole control. Polarized traps provide a temporary write torque
\cite{lapointe2011}. Shaped interfaces offer a mechanical alternative because
deviatoric curvature couples directly to a capillary quadrupole
\cite{cavallaro2011}. These measurements test the predicted separation of
short spatial correlations from long temporal inheritance.
A pair calibration is essential because \(J\) is not fixed by the far-field
quadrupole. It may contain steric and near-field capillary harmonics. Measuring
\(p(\theta_i,\theta_j\mid r)\) before the collective experiment separates the
two terms in Eq.~\eqref{eq:model}. The same microscope can then verify that
the neighbour graph remains stable throughout the read interval.

Two orthogonal polarization patterns give a direct sequence experiment. The
\(AB\) and \(BA\) histories receive the same duration and optical dose, but in
opposite order. Reading after both fields vanish tests whether the layer acts
as a passive last-event sensor. A transient surface-curvature pattern provides
a nonoptical write channel; flattening the interface removes that torque
before readout.

The result is a finite-time orientational-memory regime on a quenched
positional graph. Equilibrium glass order and operation in a mobile monolayer
are distinct questions. They reduce to measurable tests: time-step
convergence, retention under translational Brownian motion, and
independent-noise decoding of the two-pulse protocol. The Supplemental
Material gives the protocols and predeclared criteria.

Bond geometry can therefore store preparation history without selecting a
common director. Local incompatibility creates alternative angular basins,
and interactions delay escape from the selected basin. This mechanism turns
positional disorder from a source of orientational noise into a carrier of
rewritable material memory.

\begin{acknowledgments}
The mechanism hypothesis was developed with
\href{https://synthetix.institute}{Hyperion} (research build, 23 August 2026),
a domain-specific AI system that proposes physical mechanisms from equation
structure. Hyperion generated the bond-frame candidate and discriminating
tests. OpenAI Codex (GPT-5, August 2026) assisted with simulation, plotting,
code inspection, and numerical audits. The author selected the model,
interpreted every result, verified the code and claims, and takes full
responsibility for the work. Code is versioned at
\url{https://github.com/vbaulin/geometry-encoded-orientational-memory}; the
private review repository will be made public at submission. Raw records and
derived reports will be deposited on Zenodo, and the DOI will replace this
statement before submission.
\end{acknowledgments}

\begin{thebibliography}{23}

\bibitem{keim2019}
N. C. Keim, J. D. Paulsen, Z. Zeravcic, S. Sastry, and S. R. Nagel,
Memory formation in matter,
Rev. Mod. Phys. \textbf{91}, 035002 (2019).

\bibitem{keim2011}
N. C. Keim and S. R. Nagel,
Generic transient memory formation in disordered systems with noise,
Phys. Rev. Lett. \textbf{107}, 010603 (2011).

\bibitem{paulsen2014}
J. D. Paulsen, N. C. Keim, and S. R. Nagel,
Multiple transient memories in experiments on sheared non-Brownian suspensions,
Nat. Commun. \textbf{5}, 3182 (2014).

\bibitem{rotunno2005}
M. Rotunno, T. Bellini, Y. Lansac, and M. A. Glaser,
Nematics with quenched disorder: Pinning out the origin of memory,
Phys. Rev. Lett. \textbf{94}, 097802 (2005).

\bibitem{zheng2011}
Z. Zheng, F. Wang, and Y. Han,
Glass transitions in quasi-two-dimensional suspensions of colloidal ellipsoids,
Phys. Rev. Lett. \textbf{107}, 065702 (2011).

\bibitem{roller2021}
J. Roller, A. Laganapan, J.-M. Meijer, M. Fuchs, and A. Zumbusch,
Observation of liquid glass in suspensions of ellipsoidal colloids,
Proc. Natl. Acad. Sci. USA \textbf{118}, e2018072118 (2021).

\bibitem{loudet2005}
J.-C. Loudet, A. M. Alsayed, J. Zhang, and A. G. Yodh,
Capillary interactions between anisotropic colloidal particles,
Phys. Rev. Lett. \textbf{94}, 018301 (2005).

\bibitem{luo2019}
A. M. Luo, J. Vermant, P. Ilg, Z. Zhang, and L. M. C. Sagis,
Self-assembly of ellipsoidal particles at fluid--fluid interfaces with an empirical pair potential,
J. Colloid Interface Sci. \textbf{534}, 205 (2019).

\bibitem{besseling2026}
T. H. Besseling \textit{et al.},
An equilibrium rotator glass-forming phase for long-ranged repulsive colloidal rods,
Nat. Commun. \textbf{17}, 2410 (2026).

\bibitem{melio2026}
J. Melio, M. van Hecke, S. E. Henkes, \textit{et al.},
Pivoting colloidal assemblies exhibit mechanical metamaterial behaviour,
Nature \textbf{651}, 632 (2026).

\bibitem{trevenen2023}
S. Trevenen, M. A. Rahman, H. S. C. Hamilton, A. E. Ribbe,
L. C. Bradley, and P. J. Beltramo,
Nanoscale porosity in microellipsoids cloaks interparticle capillary attraction at fluid interfaces,
ACS Nano \textbf{17}, 11892 (2023).

\bibitem{rahman2026}
M. A. Rahman and P. J. Beltramo,
Controlling self-assembly and interfacial mechanics of polymer spheres and ellipsoids at fluid interfaces with surface roughness,
Soft Matter \textbf{22}, 315 (2026).

\bibitem{cavallaro2011}
M. Cavallaro, Jr., L. Botto, E. P. Lewandowski, M. Wang, and K. J. Stebe,
Curvature-driven capillary migration and assembly of rod-like particles,
Proc. Natl. Acad. Sci. USA \textbf{108}, 20923 (2011).

\bibitem{caso2025}
M. J. Caso \textit{et al.},
Capillary wave-assisted colloidal assembly,
Langmuir \textbf{41}, 3033 (2025).

\bibitem{huerre2018}
A. Huerre, M. De Corato, and V. Garbin,
Dynamic capillary assembly of colloids at interfaces with 10,000\(g\)
accelerations,
Nat. Commun. \textbf{9}, 3620 (2018).

\bibitem{lapointe2011}
C. P. Lapointe, T. G. Mason, and I. I. Smalyukh,
Towards total photonic control of complex-shaped colloids by vortex beams,
Opt. Express \textbf{19}, 18182 (2011).

\bibitem{wu2012}
H. Wu, W. Hu, H.-C. Hu, X.-W. Lin, G. Zhu, J.-W. Choi,
V. Chigrinov, and Y.-Q. Lu,
Arbitrary photo-patterning in liquid crystal alignments using a DMD-based
lithography system,
Opt. Express \textbf{20}, 16684 (2012).

\end{thebibliography}

\end{document}
